\chapter{Conclusion}
\label{ch:conclusion}

This proposal has presented the engineering architecture and evaluation plan for \textbf{Dawa Lens}, an offline-first smart medicine reminder application engineered specifically to resolve the acute challenges of medication non-adherence and pharmacological safety in Uganda's outpatient healthcare environment.

The research is motivated by the pervasive public health consequences of medication mismanagement, exacerbated by unlabelled generic pharmaceutical packaging, total omission of indigenous dietary contraindications (\textit{Matooke} and \textit{Mukene}) from standard clinical databases, frequent electrical and cellular network outages, and the absence of integrated tools to verify licensed community pharmacies. Through an exhaustive review of literature and existing commercial systems, it has been established that current technologies fail in low-resource settings because they remain tethered to continuous cloud connectivity, ignore local nutritional realities, and impose unsustainable recurring subscription or hardware costs on impoverished households.

The proposed system addresses these systemic gaps through an offline-first, four-tier architecture combining on-device optical character recognition for packaging digitization, localized biochemical interaction screening, deterministic local notification scheduling, and multi-profile family caregiver synchronization. Supported by a responsive web extension for longitudinal clinical reporting, Dawa Lens transforms any standard smartphone into an intelligent, pocket-sized clinical companion that operates with zero network dependency.

The methodology formulated in Chapter~\ref{ch:methodology} provides a disciplined, structured path grounded in Design Science Research, integrating iterative Agile sprint cycles, structured unit and integration testing, and an empirical field usability evaluation across 50 participants using the standardized System Usability Scale. The proposed project is fully achievable within the allocated academic timeframe and available computing resources at Soroti University.

Upon completion, Dawa Lens is expected to deliver a measurable $\ge 35\%$ increase in scheduled dose adherence among chronic disease outpatients, eliminate preventable adverse reactions arising from indigenous dietary interactions, provide multi-generational families with shared adherence oversight, and contribute novel empirical insights to the field of mobile health engineering in Sub-Saharan Africa. The student is firmly committed to upholding the ethical responsibilities and clinical safety standards detailed in this proposal and to working under the guidance of the assigned academic supervisor, Dr. Musa Chemisto, to deliver a high-quality, production-ready final-year engineering research project.